\chapter{StrataOS}\label{ch:os}

StrataOS 0.1 is the operating system of the prototype platform: a microkernel with preemptive
multitasking, written in \IS assembly language. It exists to show that the architecture carries a
complete system, and it uses every mechanism of Part~\ref{part:arch} for the purpose it was designed
for: user mode and the core window to confine programs, revocation by address to reclaim their memory,
endpoints for messages, core control for preemption and doorbells, the timer's WAIT register for the
time slice, \insn{svc} to block, the user-mode inbox rule to demultiplex messages, and code paging for
programs and for the kernel itself.

Its design follows the architecture's first principle. On \IS a program's working memory is its Local
Stores, and they belong to the core, not to the program. A context switch therefore has to \emph{move}
the program's Local Store contents out to far memory and another program's in, with transfers, and
StrataOS makes that cost explicit: it moves only what the program uses, it counts every byte, and it
shows the count.

\section{Structure}\label{sec:os-shape}

\begin{figure}[htbp]
\centering
\begin{tikzpicture}[
  core/.style={draw=accent,line width=0.8pt,rounded corners=3pt,minimum width=29mm,minimum height=36mm},
  k/.style={draw=accent,fill=accentlight,minimum width=25mm,minimum height=7mm,font=\footnotesize,align=center},
  p/.style={draw=black!60,fill=black!5,minimum width=25mm,minimum height=6mm,font=\footnotesize,align=center},
  >={Stealth[length=2mm]}]
  \foreach \i/\x/\t in {0/0/Core 0,1/3.3/Core 1,2/6.6/Core 3,3/9.9/Core 6} {
    \node[core] (c\i) at (\x,0) {};
    \node[font=\small\bfseries,anchor=north] at (c\i.north) {\t};
  }
  \node[k,minimum height=22mm] at (0,-0.35) {system core\\[2pt]processes\\memory\\timer, slices};
  \node[p] at (3.3,0.55) {\texttt{proc} (server)};
  \node[p] at (3.3,-0.2) {\texttt{sh}};
  \node[k] at (3.3,-1.15) {dispatcher};
  \node[p] at (6.6,0.55) {\texttt{fs} (server)};
  \node[p] at (6.6,-0.2) {\texttt{cat}};
  \node[k] at (6.6,-1.15) {dispatcher};
  \node[p] at (9.9,0.55) {\texttt{mandel} thread 0};
  \node[p] at (9.9,-0.2) {\texttt{spin}};
  \node[k] at (9.9,-1.15) {dispatcher};
  \node at (11.9,0) {\Large$\cdots$};
  \draw[->,accent,thick,dashed] (0.6,1.95) to[bend left=18] node[above,font=\scriptsize]{stop requests: time slice, doorbells} (9.9,1.85);
  \draw[<->,thick] (3.9,-2.0) to[bend right=25] node[below,font=\scriptsize]{messages: lent buffer capabilities through endpoints} (6.0,-2.0);
\end{tikzpicture}
\caption{StrataOS: the system core, and a dispatcher under the processes of every worker core}
\label{fig:os}
\end{figure}

\begin{description}[leftmargin=1em,style=nextline]
\item[The system core] Core 0 runs only kernel code, in supervisor mode, and no processes. It does
  what nothing else can: it creates and destroys processes and places them on cores, hands out memory
  (\term{arenas}, 1~MiB by default, revoked by address when a process ends), owns the timer and its single DEADLINE
  register, and drives preemption and the timeouts of timed waits. Its message interface is a handful
  of operations (Table~\ref{tab:os-kops}), and only two servers may use it.
\item[The dispatchers] Every other core is a \term{worker}. The kernel code on a worker is its
  \term{dispatcher}: the supervisor that runs the core's processes. It keeps the core's run queue,
  performs context switches, takes every capability that arrives in the core's inbox and hands it to
  the process it is for, serves \insn{svc}, emulates the inbox instructions, and pages programs' code
  units.
\item[Servers] Everything else is a user-mode \term{server} process, scheduled like any other and
  pinned to one core, at most one server per core (Table~\ref{tab:os-servers}). A request addressed to a
  core goes to that core's server.
\item[Programs] Ordinary processes: the shell, the utilities, the examples, the assembler.
\end{description}

\begin{table}[htbp]
\centering\small
\caption{The servers}\label{tab:os-servers}
\begin{tabularx}{\linewidth}{@{}>{\ttfamily}l X@{}}
\toprule
\textrm{\bfseries Server} & \textbf{Owns and does}\\\midrule
proc & the process server (pid 1): spawning, the process tree, exit statuses, WAIT and KILL,
  interrupts, the name registry. It is the endpoint every process is born with, and it holds the
  kernel's endpoint\\
dev & the device manager: holds the device root, describes the platform, hands out device
  capabilities (the disk to \texttt{fs}, the console to \texttt{tty}, the framebuffer to \texttt{ws}),
  halts the machine\\
fs & the file server: owns the block device; files, their tags, program images\\
tty & the console server: owns console output, input, break and release; line discipline, the
  foreground process, Ctrl-C\\
pipe & pipes\\
ws & the window server, started on demand: owns the framebuffer and, while it runs, the keyboard
  (Section~\ref{sec:os-ws})\\
\bottomrule
\end{tabularx}
\end{table}

The kernel shares no memory with processes and almost none between its own cores. The system core and
the dispatchers coordinate through the \term{kernel shared table}, 32~KiB of system memory that every
core reaches through a copy of one \emph{atomic} capability: the system core posts job records, kills
and notifications there, and each dispatcher publishes its run-queue length, its counters and its exit
events. All access is by far atomics, which is the only way several cores can touch the same memory.

\section{Memory}\label{sec:os-mem}

\begin{table}[htbp]
\centering\small
\caption{Local Store layout of a worker core under StrataOS}\label{tab:os-ls}
\begin{tabularx}{\linewidth}{@{}l l X@{}}
\toprule
\textbf{Store} & \textbf{Range} & \textbf{Contents}\\\midrule
I-LS & \hex{0000}--\hex{3FFF} & the boot ROM\\
I-LS & \hex{4000}--\hex{7FFF} & the resident kernel (the dispatcher)\\
I-LS & \hex{8000}--\hex{FFFF} & the running program: its code, then its code cache\\
D-LS & \hex{0000}--\hex{9FFF} & the running program's data\\
D-LS & \hex{A000}--\hex{DFFF} & kernel data and the kernel stack\\
D-LS & \hex{E000}--\hex{FFFF} & the running program's stack\\
\bottomrule
\end{tabularx}
\end{table}

The system core uses the program ranges for itself: its I-LS \hex{8000}--\hex{FFFF} is the kernel's code
cache (Section~\ref{sec:os-kpaging}) and its D-LS \hex{0000}--\hex{9FFF} holds the kernel's tables.

System memory is used as follows: $[0, 32~\text{KiB})$ the kernel shared table; $[32, 64~\text{KiB})$ the
system core's own message buffers; the \term{arena pool}, $[1, 15~\text{MiB})$ and $[16, 48~\text{MiB})$
(46~MiB; the boot ROM's load area lies between); and the framebuffer window $[48, 64~\text{MiB})$, which the
kernel makes a revocable region of its own at boot and lends to the window server.

\subsection{Arenas}
A process's memory is one arena: a run of whole 256~KiB granules of the arena pool (1~MiB unless the
spawner asked for more, Section~\ref{sec:os-servers}), zeroed when it is handed out. At boot the kernel
makes each of the pool's two parts one revocable region (\insn{cdom}) and keeps only their revokers, in
the system core's save table, and a bitmap of the pool's granules in its D-LS. It allocates an arena
from the bitmap (the first free run that is long enough) and makes a capability for it by address: it
cuts the revoker at the arena's bounds, revokes that piece with \insn{crevoke}, which kills every
capability an earlier owner of that memory may still have and returns a fresh one, and joins the
revoker again (Section~\ref{sec:domains}). When the process ends, for whatever reason, the kernel
revokes the arena the same way: every capability the process held, derived or lent out (message
buffers held by servers, window-server bands, block-device buffers), in whatever register, save table
or inbox of whatever core, dies at that instant, and the granules are free for the next process. The
number of processes is therefore limited by memory, and by the fixed tables of
Table~\ref{tab:os-limits}.

An arena larger than 1~MiB keeps the layout of Table~\ref{tab:os-arena} at its top: the offsets there
are from its \term{layout base}, $E$ bytes above its start for an arena of $1~\text{MiB} + E$, and the
program's own memory (\reg{c0}) is $E$ bytes larger, $[-E, P)$. The start-up block stays at offset
\hex{7F000} of \reg{c0}.

\begin{table}[htbp]
\centering\small
\caption{Layout of an arena for a process with $N$ threads; $P = \hex{D0000} - (N-1) \cdot 48$~KiB}\label{tab:os-arena}
\begin{tabularx}{\linewidth}{@{}l X@{}}
\toprule
\textbf{Offset} & \textbf{Contents}\\\midrule
$[0, P)$ & the program's own memory: thread 0 receives it in \reg{c0}. The image's \texttt{.far}
  section is loaded here at its own offsets, so \texttt{.far} labels are arena offsets\\
\hex{7F000} & the start-up block written by \texttt{proc}: pid, parent, the three standard handles,
  the command line\\
$[P, \hex{D0000})$ & the save areas of threads $N-1 \ldots 1$, 48~KiB each\\
$[\hex{D0000}, \hex{F0000})$ & the program image as \texttt{fs} loaded it (at most 128~KiB); the dispatcher
  of each thread keeps a read-only share and reloads code from it\\
$[\hex{F0000}, \hex{FC000})$ & thread 0's save area. With one thread, the image and the save area are
  one capability, $[\hex{D0000}, \hex{FC000})$, which the dispatcher keeps in one save-table slot\\
$[\hex{FF000}, \hex{100000})$ & the message buffer (\reg{c5})\\
\bottomrule
\end{tabularx}
\end{table}

A \term{save area} holds the part of the D-LS a thread uses while the thread is switched out: its data
extent from offset 0 (at most 40~KiB) and its live stack at offset \hex{A000} (at most 8~KiB).

\section{Processes and threads}\label{sec:os-proc}

Every process time-shares worker cores; no process owns a core. A process has one or more
\term{threads}, each on a different core.

\begin{description}[leftmargin=1em,style=nextline]
\item[StrataOS programs] are written against the program library (Section~\ref{sec:os-lib}). The
  second word of their code is the marker \texttt{"SYS2"}; the kernel recognises it and starts the
  program as one thread, without a core window, with the capabilities of Table~\ref{tab:os-start}.
\item[Bare-metal programs] have no marker. They are programs written for the bare machine, such as the
  examples of Chapter~\ref{ch:examples}, and they run unchanged: the kernel gives each thread the core
  window \csr{corebase}/\csr{corecount} (Section~\ref{sec:supervised}), so the program sees itself as
  cores $0 \ldots N-1$ of an $N$-core machine, with the console in \reg{c2} and, on its core 0, its
  memory in \reg{c0}.
\item[Threads] The shell's \texttt{-n N prog} runs a program as $N$ threads on $N$ contiguous worker
  cores that host no server. Every thread runs the same image, as every core of the bare machine does.
  Each is scheduled on its own core like any process; there is no gang scheduling. Threads give
  capabilities to each other directly with \insn{cgive}, through the window.
\end{description}

\begin{table}[htbp]
\centering\small
\caption{What a StrataOS program starts with}\label{tab:os-start}
\begin{tabularx}{\linewidth}{@{}l X@{}}
\toprule
\textbf{Register} & \textbf{Contents}\\\midrule
\reg{c0} & its arena's program memory, \perm{RWX}, revocable (it dies when the process ends)\\
\reg{c1} & a send-only endpoint to \texttt{proc}\\
\reg{c2}, \reg{c3}, \reg{c4} & send-only endpoints to \texttt{fs}, \texttt{tty} and \texttt{pipe} (absent in a
  process spawned without server endpoints)\\
\reg{c5} & its 4~KiB message buffer\\
\reg{c6} & empty: capabilities that arrive with a reply land here\\
\reg{c7} & the clock: a read-only device capability for the timer's NOW and FREQ registers\\
\reg{pc}, \reg{sp} & the image's entry at I-LS \hex{8000}; the program sets \reg{sp} to \hex{10000}\\
\bottomrule
\end{tabularx}
\end{table}

A process holds no receive right, no core-control capability, no device capability but the clock, and
no memory but its arena. Everything else it obtains by asking a server, which answers with a
capability if it agrees.

\subsection{Placement and limits}
The system core places threads when a process is created. A server goes on a core that has none.
A single-threaded process goes on the least loaded worker core, preferring cores that host a server;
an $N$-thread process goes on the run of $N$ contiguous server-free cores whose busiest core is least
loaded. The shell's \texttt{-c N} pins a process to core $N$.

\begin{table}[htbp]
\centering\small
\caption{StrataOS limits}\label{tab:os-limits}
\begin{tabularx}{\linewidth}{@{}X l@{}}
\toprule
\textbf{Resource} & \textbf{Limit}\\\midrule
Threads resident on one worker core & 16, each with 10 save-table slots reserved (with a 64-slot table: 11, with 4, the server 5)\\
Threads of one process & 6\\
Processes & memory: one arena each, 1~MiB by default, from a 46~MiB pool; 64 at most (kernel, \texttt{proc}, \texttt{fs} and \texttt{ws} tables)\\
Arenas & 64 (the kernel's arena slots); each 1--255~MiB, whole 256~KiB granules\\
Program image; code and code cache; data; stack & 128~KiB; 32~KiB; 40~KiB; 8~KiB\\
Handles per process & 8\\
Message data & 4032 bytes\\
Requests queued at a server's core; requests parked by a server & 16; 16\\
Capabilities queued at a core as attachments and gifts between threads & 16\\
Pipes; bytes buffered per pipe & 8; 2~KiB\\
Files in the directory; tags on the disk & 128; 64\\
\bottomrule
\end{tabularx}
\end{table}

With 15 cores memory binds first: \texttt{tools/test-procs.sh} runs 35 processes at once. With 6
cores, the fewest StrataOS runs on, proc and tty have a core each and the three other server cores
take 16 threads each: 50 processes, and memory (46 one-MiB arenas) binds first again;
\texttt{tools/test-procs.sh} runs 43 at once there. With a 64-slot save table a core takes 11 threads
(35 processes on 6 cores), and with three threads per core, as before the save-table pool, 6 cores
held 11 processes. Before revocation by address there were 15
revocation domains, one per arena, so at most 14 processes.

\subsection{The save-table pool}\label{sec:os-pool}

Capabilities never live in memory, so a switched-out thread's \reg{c0}--\reg{c7} stay on its core, in
the capability save table: 256 slots on the default platform ($\textit{CTAB}$, from the platform
table), reached with \insn{csavex} and \insn{crestorex}. A dispatcher keeps 7 slots for its own
capabilities (the shared table, core control, the clock, the console, the device root and the two
registers it borrows from the running thread), at 57--63 where the immediate forms reach them; the
other 249 are a \term{pool} shared by the threads and the queues, with a bitmap of the slots taken. A
thread takes slots only for what cannot be rebuilt:
\begin{itemize}
\item memory: its memory \reg{c0}, its message buffer \reg{c5}, and its image and save area, which
  are adjacent in its arena and for a thread alone in its process are one capability (a thread of a
  multi-threaded process has a read-only share of the image and a save area of its own);
\item not an endpoint: it is fully described by its node and permissions, and the dispatcher mints
  it again with \insn{cmkep} and \insn{crestrict} on switch-in;
\item not a device capability (the clock in \reg{c7}, a console port, the device root): the dispatcher
  cuts it from the device root again (\insn{csplit} twice, \insn{crestrict}). Device capabilities and
  endpoints are never revocable, so the rebuilt capability is the original: the same space, base,
  length and permissions, linear or copyable as it was.
\end{itemize}
A typical thread therefore holds 3 slots while it is switched out and 1 while it runs. Every thread
record has 10 slots reserved, all a thread could ever need (a save area, an image share and all eight
registers), and a thread's slots are charged to its reservation, so they never fail: a free thread
record on a core means free slots too, and the system core places a thread only on a core with a free
record (and a free job record); CREATE fails with $-5$ otherwise. $16 \times 10 = 160$; the other 89
are the \term{general part}, for the queue entries (requests, parked requests, attachments and gifts:
48 at most), so nothing ever runs out.

The kernel sizes all of this from the platform table at boot. With a 64-slot table (the emulator's
\texttt{-S 64}) the pool is 57 slots and a core takes 11 threads with 4 slots reserved each (the
server's record 5); the 12 general slots then hold the queue entries and the extra slots of a thread
that holds more than 4 memory capabilities (at most 8, so the queues always have 4). A full general
part refuses a request (it goes back with result $-1$, like a full queue), a PARK ($-1$), drops an
attachment, and keeps a thread whose extra capability finds no slot loaded until a slot frees up.
\texttt{cores} shows the free slots of every core.

\section{Scheduling and preemption}\label{sec:os-preempt}

A dispatcher runs its core's runnable threads round robin. It gets control only by a trap: an
\insn{svc}, a privileged instruction to emulate, a code-unit miss, a fault, a \insn{halt}, or a
\cause{stop}. On every trap it first saves the running thread's integer registers, then takes
everything from the core's inbox, applies the system core's commands, completes the waits that can
complete, and picks the thread to run. If that is the thread already loaded, no switch happens.

\begin{description}[leftmargin=1em,style=nextline]
\item[The time slice] The system core keeps a read of the timer's WAIT register in flight and waits
  with one \insn{tany} for it or for its inbox. Each time slice it reads the shared table and writes a
  stop request, through core control, to every worker that has more than one runnable thread
  (Section~\ref{sec:corectl}). The worker traps with \cause{stop} and its dispatcher switches to the next
  thread. A core with a single runnable thread is never interrupted: it pays nothing for multitasking.
\item[The slice is configurable] from 100~$\mu$s to 100~ms; the default is 2~ms. The \texttt{sched}
  program shows it together with the machine's context switches per second and KiB moved per second,
  and sets it (\texttt{sched 500us}, \texttt{sched 5ms}), so the trade between responsiveness and
  switching cost can be measured rather than assumed.
\item[Doorbells] After a process sends a message it executes \insn{svc}; its dispatcher then writes a
  stop request to the destination core, so the destination's dispatcher takes the message from its
  inbox within microseconds instead of at its next trap. A stop request that finds a core in supervisor
  mode stays pending and is taken on its return to user mode, where the dispatcher discards it if the
  work is already done.
\item[Idle] A dispatcher with nothing runnable waits in \insn{tany} for its inbox. An idle core
  executes no instructions.
\item[Timed waits] A thread that waits with a timeout registers its deadline in the shared table. The
  system core programs DEADLINE with the earliest of the next slice and every core's earliest deadline,
  and rings the core whose deadline has passed.
\end{description}

The kernel is never preempted: a stop request waits while a core is in supervisor mode. Every kernel
path is short and bounded.

\section{Context switches: the cost is explicit}\label{sec:os-switch}

A thread's state is its registers, its capability registers, its CSRs, and \emph{the parts of the
Local Stores it uses}. Table~\ref{tab:os-switch} lists what a switch does with each.

\begin{table}[htbp]
\centering\small
\caption{What a context switch saves and restores}\label{tab:os-switch}
\begin{tabularx}{\linewidth}{@{}l X@{}}
\toprule
\textbf{State} & \textbf{On a switch}\\\midrule
In-flight transfers & drained first: the dispatcher waits for every tag but its own\\
\reg{x1}--\reg{x31} & saved to the thread's record in kernel data on every trap; restored before \insn{eret}\\
CSRs & \csr{epc} of the trap, the user trap registers \csr{utvec}, \csr{uepc}, \csr{ucause}, \csr{utval},
  and \csr{xslot}, \csr{ccr}, \csr{scratch}, \csr{ttag}, \csr{tdata}, \csr{tgeom}: saved to and restored
  from the record. The core window is set from the record\\
\reg{c0}--\reg{c7} & endpoints and device capabilities described (node, or base and length, and
  permissions) and rebuilt on switch-in; memory capabilities moved to slots of the save-table pool
  (Section~\ref{sec:os-pool}) and back\\
\reg{c8}--\reg{c15} & not preserved across a switch. They are the thread's own between switches (the
  dispatcher borrows \reg{c8} and \reg{c9} during a trap and returns them)\\
D-LS data & the thread's \term{data extent}, from address 0 to the end of its data section, is copied
  to its save area with \insn{tput} and back with \insn{tget}; a bare-metal program, which may use any
  address, has the extent \hex{0000}--\hex{9FFF}. A declared scratch range is skipped\\
D-LS stack & only the live part, from \reg{sp} to \hex{10000}\\
I-LS code & never saved: it is read-only. It is reloaded with \insn{tgetx} from the image in the
  thread's arena, or from the program's code mirror\\
Code cache, overlay map & not saved: the \insn{tgetx} that reloads the code invalidates the map, and
  units are paged in again on demand\\
Software cache & its CSRs are restored; its tag table and lines are data and travel with the data extent\\
\bottomrule
\end{tabularx}
\end{table}

The I-LS cannot be read back and the D-LS has no address translation, so there is nothing to remap:
every byte of a switch is a byte moved by a transfer. A switch between two small programs moves a few
kilobytes; one between two programs that each use their whole data extent moves more than a hundred.
The dispatcher counts context switches and bytes moved per process and per core; \texttt{ps},
\texttt{cores}, \texttt{uptime} and \texttt{sched} print them, and the window server draws them.

Two supervisor calls let a program reduce or correct what a switch does.

\begin{description}[leftmargin=1em,style=nextline]
\item[Scratch ranges (SET\_SCRATCH)] A program declares one range of its D-LS as scratch. The dispatcher
  leaves it out of the saved and restored extent, and its contents are undefined after the thread has
  been switched out. Buffers that are rebuilt before every use belong there: the window server's strip
  buffers and display list are scratch, which cuts its switch cost by about two thirds.
\item[Code mirrors (SET\_MIRROR)] A program that generates code installs it in its I-LS at run time
  with \insn{tgetx} from its arena. The image no longer describes the I-LS of such a program, and the
  dispatcher cannot read the I-LS to save it. The program therefore declares a \term{code mirror}: a
  range of its arena that holds a copy of its I-LS program range from \hex{8000}. From then on the
  dispatcher reloads the code from the mirror on every switch-in. Keeping the mirror current is the
  program's job, and costs nothing: the code is written to far memory before it is installed anyway.
\end{description}

\section{Emulating the inbox}\label{sec:os-emul}

In user mode \insn{ctake}, \insn{cpoll} and \insn{tany} with the inbox bit trap with \cause{priv}
(Section~\ref{sec:userinbox}). The dispatcher has already taken everything out of the inbox; a gift
from a core inside a local process's window is a gift between threads, and waits in the dispatcher's
queue tagged with its sender. The handler decodes the trapping instruction and performs it on that
queue:
\begin{itemize}
\item \insn{ctake}: if a gift from the named sender (or from any sender of the window, lowest first)
  is queued, it goes to the destination register, \csr{takefrom} semantics included, and the thread
  continues after the instruction. Otherwise the thread blocks, another thread runs, and the
  instruction is retried when a gift arrives.
\item \insn{cpoll}: the mask of window senders with a queued gift.
\item \insn{tany}: the thread blocks until a selected tag is idle or a gift is queued.
\end{itemize}
A thread blocked this way on a core whose other threads keep it busy is found by the system core,
which stops such a core at most every 10~ms so that its dispatcher looks at its inbox. The per-sender
order of gifts is preserved, so a program that is deterministic on the bare machine
(Section~\ref{sec:determinism}) computes the same results as time-shared threads.

\section{Messages}\label{sec:os-msg}

A \term{message} is a lent buffer capability. A process owns a 4~KiB message buffer, the last 4~KiB
of its arena, in \reg{c5}. To make a request it writes a header and data into the buffer, gives the
buffer to the server, and blocks; the server reads the request, writes the result into the same buffer
and gives it back. No data is copied between address spaces, because there are none: for the duration
of the request the server owns the buffer, and the client, having given it away, cannot touch it.

\begin{table}[htbp]
\centering\small
\caption{Message header (64 bytes, little-endian doublewords); data follows at offset 64}\label{tab:os-hdr}
\begin{tabular}{@{}r l l@{}}
\toprule
\textbf{Offset} & \textbf{Field} & \textbf{Contents}\\\midrule
0 & op & the operation number\\
8 & handle & a handle, a process id, or an operation-specific value\\
16 & length & bytes of data, or an operation-specific value\\
24 & mode & flags\\
32 & result & written by the server: $\ge 0$ success, $< 0$ an error (Table~\ref{tab:os-err})\\
40 & sender & the sender's process id, stamped by the receiving core's dispatcher; servers trust it\\
48 & --- & reserved\\
\bottomrule
\end{tabular}
\end{table}

\begin{table}[htbp]
\centering\small
\caption{Results common to all servers}\label{tab:os-err}
\begin{tabular}{@{}r l @{\hspace{3em}} r l@{}}
\toprule
\textbf{Result} & \textbf{Meaning} & \textbf{Result} & \textbf{Meaning}\\\midrule
$-1$ & failed, or invalid request & $-5$ & no free core\\
$-2$ & interrupted & $-6$ & no memory\\
$-3$ & would block & $-7$ & busy\\
$-4$ & no such file, program or name & $-8$ & timed out\\
\bottomrule
\end{tabular}
\end{table}

\begin{description}[leftmargin=1em,style=nextline]
\item[Sending is a user instruction.] \insn{csend c5, cE} through a send-only endpoint. Endpoints name
  cores; the kernel mints them for the cores where servers run.
\item[Receiving is the kernel's half.] The process executes \insn{svc} WAIT\_REPLY. When the buffer comes
  back to the core's inbox the dispatcher recognises it by its base and puts it in \reg{c5} again.
\item[Attachments.] A capability sent just before a message, with a second \insn{csend}, travels with
  it: the dispatcher at the destination holds it and delivers it with the next message from the same
  sending core, in \reg{c7} of a server (or \reg{c6}, then \reg{c7}, with a reply). This is how an arena
  is lent to \texttt{fs} to load an image into, and how a server returns an endpoint or a device
  capability.
\item[Waiting on several things.] A process blocks in one \insn{svc} at a time. Servers hold a request
  that cannot be answered yet (a read from an empty pipe, a WAIT for a running child), so a client waits
  simply by not having its buffer back; a timeout on the wait, and the do-not-block mode of READ and
  WAIT, cover the rest. Nothing in StrataOS polls.
\end{description}

\subsection{Supervisor calls}\label{sec:os-svc}

The function number is in \reg{a0} and the result is returned in \reg{a0}. Timeouts are in ticks,
$-1$ for none.

\begin{table}[htbp]
\centering\small
\caption{StrataOS supervisor calls}\label{tab:os-svc}
\begin{tabularx}{\linewidth}{@{}r l X@{}}
\toprule
\textbf{\reg{a0}} & \textbf{Name} & \textbf{Arguments and result}\\\midrule
1 & WAIT\_REPLY & \reg{a1} = core to ring first (the destination of a message just sent), \reg{a2} =
  timeout, \reg{a3} = base of the capability awaited (0 = the own message buffer; $-1$ = nothing: a
  sleep), \reg{a4} = the register it is to arrive in. Returns 0, $-2$ interrupted or $-8$ timed out\\
2 & RECEIVE & servers. \reg{a1} = a mask of transfer tags to wait for as well, \reg{a2} = timeout.
  Returns a reply token, with the request in \reg{c6} and its attachment in \reg{c7}; or $-1$ with
  \reg{a1} = the idle tags; or $-8$\\
3 & REPLY & servers. \reg{a1} = token; \reg{a2} bit 0: send \reg{c7} with the reply, bit 2: send
  \reg{c4} first, bit 1 (\texttt{proc} and \texttt{dev} only): mint an endpoint for core \reg{a3} and send
  it. Gives the buffer in \reg{c6} back to its owner\\
4 & YIELD & let the core's other runnable threads run\\
5 & EXIT & \reg{a1} = exit code. \insn{halt} also exits, with the code in \reg{a0}\\
6 & NOTIFY\_WAIT & \reg{a2} = timeout. Blocks until a notification; returns its bits, or $-8$\\
7 & SET\_MIRROR & \reg{a1} = arena offset of the code mirror, \reg{a2} = its length (Section~\ref{sec:os-switch})\\
8 & PARK & servers. The request in \reg{c6} is kept by the kernel; returns its id\\
9 & UNPARK & servers. \reg{a1} = id; the parked request returns to \reg{c6}\\
10 & CLOCK & the clock capability is put in \reg{c7} again\\
11 & SET\_SCRATCH & \reg{a1} = D-LS address, \reg{a2} = length of the scratch range; length 0 clears it\\
12 & SEND & \reg{a1} = the register of a send endpoint (1--7), \reg{a5} = the register of an attachment
  (0--7, or $-1$ for none), \reg{a2}--\reg{a4} as WAIT\_REPLY. The dispatcher sends the attachment and
  then the buffer in \reg{c5} through the endpoint, with nothing of this core in between, rings the
  destination and waits as WAIT\_REPLY does. $-1$, with nothing sent, if a register is empty or is not a
  send endpoint\\
\bottomrule
\end{tabularx}
\end{table}

\subsection{Demultiplexing}\label{sec:os-demux}

Every capability that arrives at a worker is taken by its dispatcher, which decides whose it is by
what it is and where it came from. The rules are applied in this order.

\begin{enumerate}
\item A piece of the arena of a process being created on this core, known from the job record the
  system core posted: it is installed in the new thread.
\item A capability whose base is the message buffer of a local process: that process's reply.
\item A capability whose base a local process named in WAIT\_REPLY: it goes to that process.
\item A capability sent by a core inside the window of a local multi-threaded process: a gift between
  its threads, queued for the emulated \insn{ctake}.
\item A 4~KiB memory capability: a request for this core's server. It is queued, with any attachment
  waiting from the same sender, and the sender's process id is stamped into its header. On a core
  without a server, or when eight requests are already queued, it is given back at once with result $-1$.
\item Anything else: an attachment. It waits for the next message from the same sending core.
\end{enumerate}

A process cannot forge a reply or a request for another process, because what distinguishes them is
the identity of memory it does not hold.

\subsection{The kernel's operations}

The system core has an endpoint too. Only \texttt{proc} holds it (and \texttt{dev}, to halt the machine
and to lend the framebuffer). Tables~\ref{tab:os-kops}--\ref{tab:os-pipe} give the complete operation
numbering of StrataOS 0.1; operation numbers are per server.

\begin{table}[htbp]
\centering\small
\caption{Operations of the system core}\label{tab:os-kops}
\begin{tabularx}{\linewidth}{@{}r l X@{}}
\toprule
\textbf{Op} & \textbf{Name} & \textbf{Request and reply}\\\midrule
1 & CREATE & attachment: an arena holding the image at \hex{D0000} and the start-up block. Data: pid,
  thread count (0 = one time-shared thread), flags, a core to pin to or $-1$, the cores of the endpoints
  for \reg{c1}--\reg{c4}, \reg{a0}--\reg{a2}. Reply: 0 and the first core; the process starts. A server's
  reply carries an endpoint to its core\\
2 & DESTROY & handle = pid. The process is stopped and its arena revoked; \texttt{proc} later
  receives EXITED with the killed bit\\
3 & MEMORY & length = the size wanted (0: 1~MiB; rounded up to whole 256~KiB granules, at least 1~MiB).
  Reply: a zeroed arena with a fresh capability, as an attachment; data: its base and size; $-6$ if none\\
4 & STATS & reply: counters for the machine, every core and every process\\
5 & TUNE & length = the time slice in ticks, 100--100000; $-1$ out of range\\
6 & FREE & attachment: an arena to return (it is revoked, and its granules are free)\\
7 & EXITED & sent \emph{by} the kernel to \texttt{proc}: handle = pid, length = exit status\\
8 & HALT & stops every core\\
9 & NOTIFY & handle = pid, length = notification bits: wakes its NOTIFY\_WAIT or timed sleep\\
10 & FBLEND & \texttt{dev} only. Handle = pid: revokes every older capability of the framebuffer
  and replies with a fresh one for that process; the kernel revokes it again when the process
  ends. $-7$ if a live process holds it\\
\bottomrule
\end{tabularx}
\end{table}

\section{The servers}\label{sec:os-servers}

\begin{table}[htbp]
\centering\small
\caption{Operations of \texttt{proc} (endpoint \reg{c1})}\label{tab:os-proc}
\begin{tabularx}{\linewidth}{@{}r l X@{}}
\toprule
\textbf{Op} & \textbf{Name} & \textbf{Request and reply}\\\midrule
5 & SPAWN & data: thread count (0 or 1 = one thread, $N$ = $N$ threads on $N$ cores), three handle
  descriptors that become the child's handles 0, 1, 2 ($-1$ = none), then the command line
  \texttt{"prog args"}. Mode bit 0: no server endpoints; bit 1: pin to core (mode $\gg 8$) \& 255;
  mode $\gg 16$: the arena's size in MiB (0: 1~MiB). Reply: the pid; $-4$, $-5$, $-6$\\
6 & WAIT & handle = pid, or 0 for any child; mode bit 0 = do not block. Reply: the exit status, with
  the pid in the data; $-3$ still running; $-2$ interrupted; $-1$ no child\\
7 & KILL & handle = pid\\
9 & PS & reply: the count, and one 32-byte record per process: u8 pid, state (1 running, 2 ready,
  3 blocked, 4 zombie), core, parent, flags (2 = server), threads; 2 reserved; u32 switches, u32 KiB
  moved; 16 bytes of name\\
16 & INFO & reply data: cores, free heap, total memory, processes, kernel units loaded, context
  switches, KiB moved (u64 each)\\
17 & SIGNAL & handle = pid: interrupt it\\
22 & CATCH & handle = 1 catch interrupts, 0 do not\\
25 & LOOKUP & data = a server's name. Reply: 0, with a send-only endpoint to it as an attachment; $-4$\\
29 & REGISTER & data = a name: the caller becomes the server of that name; the entry dies with it.
  $-5$ unless the caller was started as its core's server (the boot servers, and \texttt{ws}); $-7$ if
  another live process holds the name\\
30 & SCHED & handle 0: reply data = the time slice, switches per second, KiB moved per second since
  the last read. Handle 1: length = a new slice, forwarded to TUNE\\
31 & CORES & reply: the count, and one 32-byte record per core: u8 core, pid running, runnable,
  processes, flags (2 = server); 3 reserved; u32 switches, u32 KiB moved\\
39 & WHOIS & data = a name. Reply: the pid registered under it (0 for the boot servers); $-4$ if
  none. \texttt{dev} asks it before GRANT\_FB\\
\bottomrule
\end{tabularx}
\end{table}

An exit status has the exit code in bits 31:0; bit 62 is set if the process was killed, and bit 63 if
it faulted, with the trap cause in bits 7:0.

\textbf{Spawning} shows how the pieces combine. \texttt{proc} asks the kernel for an arena (MEMORY),
writes the start-up block into it, lends it to \texttt{fs} to read the program image into (LOADIMG),
tells the servers of the three handles that the child now holds them (DUPS), and asks the kernel to
CREATE the process from the arena. The system core chooses the cores, cuts the arena into the pieces
of Table~\ref{tab:os-arena}, posts a job record for each thread and gives each worker its pieces; when
a dispatcher has them all it makes the thread runnable. While it waits for these answers, \texttt{proc}
keeps the caller's request parked. When a process ends, its dispatcher reports the exit, the system
core revokes the arena and sends EXITED, and \texttt{proc} tells the servers to drop the process's
handles (CLEANUP) and completes the parent's WAIT.

\begin{table}[htbp]
\centering\small
\caption{Operations of \texttt{dev} (through LOOKUP)}\label{tab:os-dev}
\begin{tabularx}{\linewidth}{@{}r l X@{}}
\toprule
\textbf{Op} & \textbf{Name} & \textbf{Request and reply}\\\midrule
23 & DEVS & reply: the count and the platform table's 32-byte entries (Table~\ref{tab:stpdent})\\
24 & HALT & asks the kernel to stop every core\\
26 & GRANT\_FB & reply: 0, with the framebuffer memory and a device capability for the controller's
  PRESENT and VSYNC registers as attachments. $-5$ unless the caller is the process registered as
  \texttt{ws}; $-7$ while a live process holds the framebuffer\\
30 & GETDISK & for \texttt{fs}: a device capability for the block device's window and a send-only
  endpoint to its node\\
31 & GETCON & for \texttt{tty}: a device capability for the console's ports\\
59 & GETNET & for \texttt{net}: a device capability for the network interface's ports. $-5$ unless
  the caller is the process registered as \texttt{net}\\
\bottomrule
\end{tabularx}
\end{table}

\begin{table}[htbp]
\centering\small
\caption{Operations of \texttt{fs} (endpoint \reg{c2})}\label{tab:os-fs}
\begin{tabularx}{\linewidth}{@{}r l X@{}}
\toprule
\textbf{Op} & \textbf{Name} & \textbf{Request and reply}\\\midrule
1 & READ & handle, length. Reply: bytes read, in the data; 0 at the end of the file\\
2 & WRITE & handle, length, data. Reply: bytes written\\
3 & OPEN & data = a path (Section~\ref{sec:os-tags}); mode 0 read, 1 write (truncate), 2 append. Writing
  gives the file the path's tags. Reply: a file id; $-4$\\
4 & CLOSE & handle. A file that was written is saved to disk when its last reference closes\\
10 & DIR & data = a path of tags (none: every file; \texttt{-tag} excludes). Reply: the count of
  files carrying them all, and 32-byte records: 24 bytes of name, u64 size (126 at most); the mode
  word holds the facets, a tag set of the other tags those files carry; $-4$ no such tag\\
11 & UNLINK & data = a path; $-4$\\
12 & TAG & data = a path, a zero byte, a tag name (created if new). $-4$; $-6$ all 64 tags in use;
  $-1$ a bad tag name\\
13 & UNTAG & data = a path, a zero byte, a tag name. $-4$\\
14 & TAGS & data = a path (none: every tag). Reply: the count, and 32-byte records: 16 bytes of name,
  u64 files carrying it, u64 tag number: every tag, or the tags of that file\\
15 & RENAME & data = a path, a zero byte, a new path. The file takes the new name and gains the new
  path's tags (created if new); it keeps its tags, blocks and directory entry, and another file of
  that name is deleted. $-4$; $-6$; $-1$ a bad new name\\
20 & CWD & the requester's current tag path (Section~\ref{sec:os-cwd}). Mode 0: reply = its length,
  data = its text (\texttt{/} at the root). Mode 1: data = a path, relative to the current one
  (\texttt{..} allowed; none: the root); reply as mode 0; $-4$ no such tag, or no file in that view;
  $-1$ longer than 64 bytes. Mode 2, from \texttt{proc} only: data = u64 parent pid, u64 child pid:
  the child starts in the parent's path\\
27 & STAT & data = a path. Reply: the size; $-4$\\
28 & LOADIMG & data = a path, mode = arena offset, handle = the most bytes; the arena as
  attachment. Reply: the image's size; the arena comes back with the reply\\
32 & STATFS & reply data: total blocks, used blocks, files\\
33 & READBLK & handle = a block number. Reply data: the raw 512-byte block\\
34 & READALL & handle = a file id, mode = arena offset, length = the most bytes; the arena as
  attachment: the whole file is copied into it\\
35 & WRITEALL & handle = a file id, mode = arena offset, length; the arena as attachment: appended to
  the file\\
\bottomrule
\end{tabularx}
\end{table}

\texttt{fs} drives the block device entirely from user mode (Section~\ref{sec:blkdev}): it writes a
command to the submission port through its device capability, lends the device its I/O buffer with
\insn{csend} through the endpoint to the device's node, and waits with WAIT\_REPLY for the buffer's
base. The device's return of the buffer arrives at the core's inbox like any reply. Program images are
read straight into the arena \texttt{proc} lends; the data never passes through \texttt{fs}'s Local
Store. Files are contiguous runs of blocks; a file being written is held in a 128~KiB region.

\begin{table}[htbp]
\centering\small
\caption{Operations of \texttt{tty} (endpoint \reg{c3})}\label{tab:os-tty}
\begin{tabularx}{\linewidth}{@{}r l X@{}}
\toprule
\textbf{Op} & \textbf{Name} & \textbf{Request and reply}\\\midrule
1 & READ & length; mode bit 0 = do not block. Cooked mode: input is echoed and edited with Backspace,
  and a READ returns one line with its newline, or what was typed before Ctrl-D; 0 = end of input.
  Raw mode: no echo; whatever bytes are there, at least one; a READ that does not block and finds
  nothing has one byte read ahead, so a program that polls sees the rest of an escape sequence. CR
  becomes LF in both\\
2 & WRITE & data to the console\\
4 & CLOSE & 0 (the console is never closed)\\
18 & SETFG & handle = pid: Ctrl-C now interrupts that process\\
19 & CONMODE & handle = 0 cooked, 1 raw\\
\bottomrule
\end{tabularx}
\end{table}

\texttt{tty} is the model of a driver without interrupts. It never blocks on its device: it keeps a
read of the break port in flight on one tag always, and a one-byte read of console input on another
while a reader is waiting, and its RECEIVE names both tags, so the one supervisor call wakes it for a
request, a key or Ctrl-C. Input is neither read nor echoed before a program asks for it. Because it
holds pending device reads, \texttt{tty} is alone on its core.

\begin{table}[htbp]
\centering\small
\caption{Operations of \texttt{pipe} (endpoint \reg{c4})}\label{tab:os-pipe}
\begin{tabularx}{\linewidth}{@{}r l X@{}}
\toprule
\textbf{Op} & \textbf{Name} & \textbf{Request and reply}\\\midrule
8 & PIPE & reply: read id $|$ write id $\ll 8$\\
1 & READ & handle, length; mode bit 0 = do not block. Reply: bytes; 0 = end of input (no writer is
  left). A reader of an empty pipe waits\\
2 & WRITE & handle, length, data. Reply: bytes taken, at least 1; $-1$ if no reader is left. A writer
  to a full pipe waits\\
4 & CLOSE & handle\\
\bottomrule
\end{tabularx}
\end{table}

Three further operations go from \texttt{proc} to every server that keeps handles: 36 DUPS (a process
now holds these handles too), 37 CLEANUP (a process ended: drop everything it held) and 38 CANCEL
(answer the requests held for a process with $-2$). The ends of files and pipes are reference-counted
per process, which is how a pipe's reader sees end of input when the last writer exits.

\subsection{Tags instead of directories}\label{sec:os-tags}
The file system has no directory tree. A file has one name, unique on the disk, and carries any number
of the disk's 64 tags (names of 1--16 bytes without \texttt{/}); a ``directory'' is a set of tags, so
one file shows up in several directories at once, without copies or links. Every name \texttt{fs}
takes is a \emph{path}: components separated by \texttt{/}, empty ones skipped. The last component
is the file's name and the ones before it are tags the file must carry: \texttt{motd},
\texttt{/motd} and \texttt{/doc/motd} all name \texttt{motd} when it carries \texttt{doc}, and
\texttt{/bin/motd} names nothing. Opening for writing gives the file the path's tags, creating new
ones (\texttt{echo hi > /notes/todo}), and a file that a write replaces keeps its tags. DIR lists the
files that carry every tag of its path (\texttt{ls /src/demo} is an intersection); when the last
component is not a tag but names a file that carries the others, it lists that one file. Deleting
through a path deletes the file; UNTAG takes it out of a directory. A tag no file carries any more
disappears. Programs that only open a file by name work unchanged, and can be given paths.

A component \texttt{-tag} excludes: the file must not carry that tag, so \texttt{ls /src/-demo} lists
the sources that are not demos (excluding a tag that does not exist excludes nothing). A path that
creates a file or tags (writing, RENAME's new path) cannot exclude, and no tag name starts with
\texttt{-}.

\paragraph{The current tag path.}\label{sec:os-cwd} Every process has a current tag path, the
analogue of a working directory, kept by \texttt{fs}: one entry per live process that is not at
\texttt{/} (32 entries of at most 64 bytes), found by the sender pid the kernel stamps on each
request. A path that does not start with \texttt{/} is read as that path, a \texttt{/}, then the
path; DIR with no path lists it. In \texttt{/src/demo}, \texttt{cat hello.s} opens
\texttt{/src/demo/hello.s}, \texttt{ls x/} lists \texttt{/src/demo/x}, and \texttt{echo hi > notes}
creates \texttt{notes} tagged \texttt{src} and \texttt{demo}. Every path is then normalized:
\texttt{.} components go, \texttt{..} removes the one before it. When creating, exclusions that come
from the current path are dropped (the new file simply gets its tags). Setting it (CWD mode 1, the
shell's \texttt{cd}) requires every component to be a tag that exists (or \texttt{-} and one) and
at least one file to match, so no process is left in an empty view. A child starts in its parent's
path: \texttt{proc} sends CWD mode 2 while it spawns, after LOADIMG and before CREATE, and the
CLEANUP it sends when the process ends frees the entry. \texttt{proc} itself stays at \texttt{/}, so
program images are always found from the root. Every client gets this for free: library programs,
Forth's \texttt{include} and file words, and servers.

\paragraph{Subdirectories are facets.} With its files DIR answers the other tags that at least one of
them carries: the OR of their tag sets, without the path's own tags, computed in the same loop. These
are the ways to narrow the listing further, as subdirectories would be: \texttt{ls /src} shows the
files tagged \texttt{src} and then \texttt{demo/}, \texttt{ls /src/demo} narrows to both, and
\texttt{ls /} ends with every tag in use. Nothing on the disk changes. The shell's Tab completes facets
in a path like folders (\texttt{cat /sr} gives \texttt{/src/}).

On the disk the directory keeps its 128 entries of 32 bytes (the kernel reads it at boot to load the
first servers), and the tag table follows it: four blocks, at the LBA in boot-header offset 52, of 64
tag names of 16 bytes (zero-padded; a zero first byte is a free slot), then 128 64-bit tag sets, one
per directory entry: bit $t$ of set $f$ says that entry $f$ carries tag $t$. A query is a loop over
128 words with one \texttt{and} each. \texttt{fs} keeps the table (2~KiB) in its D-LS and writes it
back with the directory. \texttt{mkdisk -t tag,tag\ldots} tags the files that follow it; the supplied
disk tags its files \texttt{doc}, \texttt{etc}, \texttt{sys}, \texttt{bin}, \texttt{ws},
\texttt{demo}, \texttt{src} and \texttt{forth}.

\subsection{Handles}
A handle is a small number, per process: 0 standard input, 1 standard output, 2 standard error, and up
to five more from OPEN and PIPE. The library keeps, for each, the server it belongs to and the
server's id for it, and sends READ, WRITE and CLOSE to that server; a program does not know whether its
output is the console, a file or a pipe. SPAWN passes three handles to the child.

\subsection{Time}\label{sec:os-time}
Time needs no call. Every process holds the clock in \reg{c7}, a read-only capability for exactly the
timer's NOW and FREQ registers (Section~\ref{sec:timer}): reading the clock is one 8-byte \insn{tget}.
Sleeping is WAIT\_REPLY with a timeout and nothing awaited. DEADLINE and WAIT belong to the system
core.

\section{The program library}\label{sec:os-lib}

Programs include \texttt{os/lib/sys.s}. It builds messages, keeps the handle table and buffers
standard output.

\begin{asm}
_start: j     main                      ; skip the marker
        .word 0x32535953                ; "SYS2": a StrataOS program
main:   li    sp, 65536                 ; the stack is at the top of the D-LS
        call  sys_init                  ; read the start-up block, set up handles 0..2
        la    a0, text                  ; ...
        li    a1, text_end - text       ; ...
        call  out_str                   ; buffered standard output
        li    a0, 0                     ; exit code
        call  sys_exit                  ; flush standard output, end the process
\end{asm}

\begin{table}[htbp]
\centering\small
\caption{The program library, \texttt{os/lib/sys.s}}\label{tab:syslib}
\begin{tabularx}{\linewidth}{@{}l >{\ttfamily\raggedright\arraybackslash}X@{}}
\toprule
\textbf{Area} & \textrm{\bfseries Routines}\\\midrule
Start and exit & sys\_init, sys\_exit(code), sys\_pid, sys\_cmdline, sys\_args(buf), arg\_next\\
Handles & sys\_read(h, buf, max), sys\_readnb, sys\_write(h, buf, len), sys\_close(h),
  sys\_open(name, len, mode), sys\_pipe, sys\_hdesc(h), sys\_hset(h, desc)\\
Processes & sys\_spawn(n, mode, cmd, len, h0, h1, h2), sys\_wait(pid, nb), sys\_kill(pid),
  sys\_signal(pid), sys\_catch(on), sys\_ps(buf), sys\_info(buf), sys\_cores(buf),
  sys\_sched(op, arg), sys\_lookup(name, len, reg), sys\_register(name, len)\\
Console & sys\_conmode(mode), sys\_setfg(pid)\\
Files & sys\_dir(buf), sys\_unlink, sys\_stat(name, len), sys\_statfs(buf), sys\_readblk(n, buf),
  sys\_rename(path, len, new, len)\\
Tags & sys\_query(buf, path, len), sys\_facets (of the last query), sys\_tags(buf, path, len),
  sys\_tag(path, len, tag, len), sys\_untag(path, len, tag, len)\\
Devices & sys\_devs(buf), sys\_halt, sys\_grantfb\\
Time & sys\_now, sys\_sleep(ticks), sys\_notify\_wait(timeout), sys\_yield, sys\_clock\\
Switch cost & sys\_mirror(offset, len), sys\_scratch(start, len)\\
Windows & ws\_call(op, handle, data, len)\\
Output and input & out\_str, out\_chr, out\_nl, out\_dec, out\_sdec, out\_decw, out\_pad, out\_hex,
  out\_flush, err\_str, in\_line(buf, max)\\
Strings & str\_atoi, str\_num, str\_eq, str\_has, str\_len\\
\bottomrule
\end{tabularx}
\end{table}

Results are returned in \reg{a0}. Calls clobber \reg{t0}--\reg{t6} and \reg{a0}--\reg{a7}; the saved
registers, \reg{sp}, \reg{gp} and \reg{tp} are kept. Transfer tag \reg{tg10} belongs to the library. A
write longer than a message is sent in pieces. Servers also include \texttt{os/lib/srv.s}: the
receive, reply and park loop, a guard that drops a request whose owner died while it was being served,
and the reference table for handles.

\subsection{Bare-metal programs and redirection}
A bare-metal program knows nothing of messages. Its console output works because it holds the console
in \reg{c2}. For redirection the shell and \texttt{proc} use a memory convention: with \texttt{<
file}, the file is copied into the arena at offset \hex{80000} (at most 256~KiB) and the program starts
with its length in \reg{a0}; with \texttt{> file}, \reg{a1} is 1, the program leaves its output at
arena offset \hex{C0000} and exits with its length (at most 64~KiB), and \texttt{proc} appends it to the
file. The self-hosted assembler uses this convention (Section~\ref{sec:os-as}).

\section{Interrupts and Ctrl-C}\label{sec:os-break}

StrataOS has one asynchronous signal, the \term{interrupt}, and builds it from three mechanisms of
Chapter~\ref{ch:traps}: an event that is a transfer, messages, and core control.

\begin{enumerate}
\item \texttt{tty} keeps a read of the console break port in flight. When the user presses Ctrl-C the
  read completes, \texttt{tty}'s RECEIVE returns with that tag idle, and \texttt{tty} echoes
  \texttt{\textasciicircum C}, drops the partial input line and asks \texttt{proc} to SIGNAL the
  \term{foreground process} (the one named by the last SETFG).
\item If the process does not catch interrupts, \texttt{proc} asks the kernel to DESTROY it. Its
  cores are stopped through core control wherever they are, its arena is revoked, and its parent's
  WAIT returns a status with the killed bit.
\item If the process catches interrupts (CATCH), the blocking call it is in, or its next one, returns
  $-2$: \texttt{proc} answers its held WAIT, tells \texttt{tty} and \texttt{pipe} to CANCEL the
  requests they hold for it, and has the kernel NOTIFY it to end a sleep or a NOTIFY\_WAIT.
\end{enumerate}

The shell catches interrupts. It gives the console to a foreground job with SETFG and takes it back
afterwards; interrupted while it waits for a job, it SIGNALs every stage of the job; interrupted while
reading a line, it abandons the line. The window server catches interrupts and forwards them to the
process of the focused window. The editor catches them too: its read returns $-2$, which it takes as
the key Ctrl-C (Section~\ref{sec:os-vi}).

\section{Code paging}

\subsection{Programs}\label{sec:os-paging}

A program can keep code that does not fit in its I-LS in its image's \texttt{.far} section, organised
in units of $2^n$ bytes, set \csr{xslot} to $n$ and call the units with \insn{xcall} through its arena
capability \reg{c0}; units that call other units return with \insn{xret ra}
(Section~\ref{sec:overlays}). Nothing is resident at first. A program that installs its own trap handler
serves its own misses, in user mode, as the \texttt{overlay} example does. Otherwise every miss traps
to the dispatcher, which:
\begin{itemize}
\item takes the far address from \csr{tval} (clearing bit~63 for a return), and rejects the miss if
  \csr{xslot} is 0, the program no longer holds \reg{c0}, \reg{c0}'s base is not a multiple of the
  slot size (\insn{xmap} needs aligned units), or the unit is not a whole slot inside it;
\item uses the I-LS from the end of the program's resident code, rounded up to a slot, to \hex{10000}
  as the code cache, at most one slot per overlay map entry ($\textit{NXMAP}$, which the kernel takes
  from the platform table), replaced round robin;
\item loads the unit into the victim slot with \insn{tgetx}, maps it with \insn{xmap} and resumes the
  program at the instruction that missed.
\end{itemize}
A miss the dispatcher cannot serve ends the program with a fault of cause \cause{xmiss}. The
\texttt{units} program runs 24 one-KiB units through sixteen slots (eight on an eight-entry platform) with call chains seven deep and
checks its result against the same computation done by resident code. After a context switch the
cache is empty and refills on demand.

\subsection{The kernel pages itself}\label{sec:os-kpaging}

The resident kernel is the 16~KiB at I-LS \hex{4000}--\hex{7FFF} of every core: the dispatcher, and on
the system core the event loop and the paging handler. Everything the system core does rarely
(start-up, creating and destroying processes, memory, statistics, exits) is \term{kernel code units}:
27 units of 2~KiB in the kernel image's \texttt{.far} section, which the boot ROM loads into the load
area and hands to core 0 as the far image share \reg{c12} (Section~\ref{sec:boot}).

Resident code calls a unit with \insn{xcall ra, c12, t5}. A miss traps to the system core's own
supervisor handler, which copies the unit into one of $\textit{NXMAP}$ cache slots (sixteen by
default) of 2~KiB from I-LS \hex{8000}, round robin, maps it and retries: the mechanism programs use, applied by the
kernel to itself, with \csr{xslot} = 11. Units run at a cache address rather than where they were
assembled, so they call resident code through a register, call other units with \insn{xcall}, and
return with \insn{xret}, whose far link pages the caller back in if it was evicted meanwhile.
\texttt{uptime} reports how many units have been paged in.

\section{The shell}\label{sec:os-sh}

\begin{shell}
sh                  interactive when standard input is the console
sh -i               interactive even if standard input is a pipe
sh -l               login shell: run the script rc first; at end of input, halt the machine
sh -c "ls | wc"     run one command line and exit with its status
sh FILE             run each line of FILE, then exit
\end{shell}

\begin{table}[htbp]
\centering\small
\caption{Shell syntax}\label{tab:shell}
\begin{tabularx}{\linewidth}{@{}>{\ttfamily}l X@{}}
\toprule
\textrm{\bfseries Form} & \textbf{Meaning}\\\midrule
a ; b & run \texttt{a}, then \texttt{b}\\
a \&\& b \quad a || b & run \texttt{b} only if \texttt{a} succeeded, or failed\\
a \& & run \texttt{a} in the background: \texttt{[J] PID} now, \texttt{[J] done  a} later\\
a | b | c | d & a pipeline of up to four stages\\
< f \quad > f \quad >> f & redirect the first stage's input, the last stage's output (write, append)\\
'text' \quad "text" & one word; \texttt{\$} is expanded inside double quotes only; \texttt{\textbackslash x} is the character \texttt{x}\\
\$? \quad \$\$ \quad \$NAME & the last status, the shell's pid, a variable\\
-n N prog & run \texttt{prog} as $N$ threads on $N$ cores\\
-c N prog & pin \texttt{prog} to core $N$\\
-m N prog & give \texttt{prog} an arena of $N$~MiB (1--255; the default is 1)\\
run prog & the same as \texttt{prog}\\
time cmd & run \texttt{cmd}, then print the wall time it took\\
\# comment & the rest of the line is ignored\\
\bottomrule
\end{tabularx}
\end{table}

The status \texttt{\$?} is the exit code, 130 for a process killed by Ctrl-C, and 128 plus the trap
cause after a fault; a non-zero status shows in the prompt, as in \texttt{[130] strata\$}.

\begin{description}[leftmargin=1em,style=nextline]
\item[Builtins] \texttt{help}, \texttt{echo}, \texttt{set}, \texttt{unset}, \texttt{history},
  \texttt{jobs}, \texttt{fg}, \texttt{wait}, \texttt{kill} (a pid, or \texttt{\%J} for every stage of job
  $J$), \texttt{sleep} (decimal seconds), \texttt{clear}, \texttt{source}, \texttt{exit}, \texttt{halt},
  \texttt{true}, \texttt{false}, \texttt{cd}, \texttt{pwd}. A builtin that is a pipeline stage or runs
  in the background is run by a child \texttt{sh -c}.
\item[cd and pwd] \texttt{cd [PATH]} sets the current tag path (Section~\ref{sec:os-cwd}):
  \texttt{cd TAG} appends a tag, \texttt{cd a/b} both, \texttt{cd -x} an exclusion, \texttt{cd /x/y}
  replaces the path, \texttt{cd ..} drops the last component, \texttt{cd} or \texttt{cd /} goes back to
  every file. A tag no file carries is refused (\texttt{sh: cd: no such tag: PATH}, status 1).
  \texttt{pwd} prints the path. Away from the root the prompt shows it, as in
  \texttt{strata:/src/demo\$}; programs and child shells start in it.
\item[Line editing] The shell switches the console to raw mode while a line is typed. Left, Right,
  Home, End and Ctrl-B/F/A/E move; Backspace and Delete erase; Ctrl-K, Ctrl-U and Ctrl-W erase to the
  end, to the start, and the word before the cursor; Up, Down and Ctrl-P/N walk a history of 32 lines;
  Ctrl-L clears the screen; Tab completes a command (builtins and the programs on the disk) or a file
  name, also behind a tag path (\texttt{cat /doc/mo}; a tag completes as \texttt{tag/}) and under the
  current tag path (its files and facets), and a second Tab lists the candidates; Ctrl-D on an empty line ends the input; Ctrl-C abandons
  the line.
\item[Running a command] Each stage's command text goes to \texttt{proc} with SPAWN, together with the
  three handles that become the child's 0, 1 and 2: the shell's own, pipe ends from \texttt{pipe}, or
  files from \texttt{fs}. The shell closes its copies after each spawn, gives the console to the last
  stage with SETFG, and waits for every stage.
\item[Where its state lives] The history and the scripts being run are in the shell's arena, reached
  with transfers; its D-LS holds only the lines being worked on, which keeps its context switches
  small.
\end{description}

At boot the login shell runs the script \texttt{rc} if the disk has one; the supplied one prints
\texttt{motd}. When console input ends, the login shell halts the machine.

\subsection{Programs on the disk}

\begin{table}[htbp]
\centering\small
\caption{Programs supplied on the StrataOS disk}\label{tab:os-progs}
\begin{tabularx}{\linewidth}{@{}>{\ttfamily}l X@{}}
\toprule
\textrm{\bfseries Program} & \textbf{Action}\\\midrule
ls [-t] [PATH] & list the files with their sizes: those of the current tag path (every file at
  \texttt{/}), or those carrying the tags of PATH
  (\texttt{ls /src/demo}, \texttt{ls /src/-demo}), then its facets as \texttt{tag/};
  \texttt{-t} adds each file's tags\\
tag FILE TAG \ldots & give a file tags; \texttt{untag FILE TAG \ldots} takes them away;
  \texttt{retag FILE [TAG \ldots]} leaves it exactly these\\
tags [FILE \ldots] & every tag with the number of files carrying it, or the tags of each FILE\\
cat [NAME \ldots] & copy the named files, or standard input, to standard output\\
rm NAME \ldots & delete files\\
cp [-t] SRC DST & copy a file; DST's path gives the copy tags, \texttt{-t} also SRC's (DST must have
  another name: names are unique)\\
mv SRC DST & rename (RENAME): the file keeps its tags and gains DST's; \texttt{mv f /todo/} keeps the
  name\\
touch FILE \ldots & create empty files; a file the path does not name yet gets the path's tags\\
df & disk blocks: total, used, free; the number of files\\
hexdump BLOCK & the first 256 bytes of a raw disk block\\
ps & the processes: pid, parent, state, core, threads, context switches, KiB moved, name\\
cores & what every core is doing: the running process, runnable and resident counts, switches and
  KiB moved\\
free & total and free memory, and the number of processes\\
uptime & time since power-on, processes, context switches and KiB moved, kernel units paged in\\
sched [TIME] & show the time slice and the switch rates, or set the slice
  (\texttt{500us}, \texttt{5ms}, or a number of microseconds)\\
devs & the platform table\\
seq N & print 1 to $N$, one per line\\
grep TEXT & print the lines of standard input that contain TEXT\\
wc & the line, word and byte counts of standard input\\
upper & copy standard input to standard output in capital letters\\
spin [N [C]] & burn CPU, printing a character now and then: $N$ times, or forever\\
units & demand-paged code (Section~\ref{sec:os-paging})\\
as & the assembler (Section~\ref{sec:os-as})\\
sh & the shell\\
vi [FILE] & the editor (Section~\ref{sec:os-vi})\\
\bottomrule
\end{tabularx}
\end{table}

The disk also carries the five servers, the bare-metal examples \texttt{hello}, \texttt{stream\_sum},
\texttt{swcache}, \texttt{overlay}, \texttt{faults}, \texttt{pipeline}, \texttt{parsum}, \texttt{mandel}
and \texttt{race}, the sources of all of these (a utility's source has the library appended, without
its comments, so that they fit the 4~MiB disk), and the text files \texttt{rc}, \texttt{motd} and
\texttt{about}. A session:

\begin{shell}
strata$ seq 30 | grep 2 | wc
12 12 35
strata$ -c 6 spin 12 a > sa & -c 6 spin 12 b > sb &
[1] 9
[2] 10
strata$ ps
strata$ -n 3 mandel
strata$ sched 500us
\end{shell}

The pipeline is three processes and two pipes; the two \texttt{spin} processes share core 6 and
\texttt{ps} shows the switches between them and the bytes those switches moved; \texttt{mandel} runs as
three threads on three cores, its \insn{ctake}s emulated.

\subsection{The editor}\label{sec:os-vi}
\texttt{vi} \textit{file} (\texttt{os/progs/vi.s}) edits a file in the manner of vi, on the console and
in a \texttt{term} window, with Strata assembly and Strata Forth in colour. It has normal, insert,
command-line (\texttt{:}) and search (\texttt{/}, \texttt{?}) modes and a status line (the mode, the
file, \texttt{[+]} when modified, line and column). It moves with \texttt{h j k l} and the arrows,
\texttt{w b e}, \texttt{0 \^{} \$}, \texttt{gg G}, \textit{N}\texttt{G}, \texttt{\%} and Ctrl-D, U, F, B;
inserts with \texttt{i a I A o O}; changes with \texttt{x X r D C s S J}, the operators \texttt{d c y}
with a motion or doubled (\texttt{dw}, \texttt{cw}, \texttt{d\$}, \texttt{dd}, \texttt{3dd}),
\texttt{Y}, \texttt{p P}, \texttt{u} (undo, many levels), Ctrl-R (redo) and \texttt{.} (the last change,
its keys replayed); searches with \texttt{/} \texttt{?} \texttt{n} \texttt{N} (plain text); and takes
\texttt{:w}, \texttt{:w} \textit{name}, \texttt{:q}, \texttt{:q!}, \texttt{:wq}, \texttt{:x},
\texttt{:e} \textit{name}, \texttt{:}\textit{N}, \texttt{:set nu}, \texttt{:set nonu},
\texttt{:syntax on} and \texttt{off}.

\begin{description}[leftmargin=1em,style=nextline]
\item[Highlighting] \texttt{.s} files are Strata assembly, \texttt{.fs} files Forth (another name: by
  the file's tags \texttt{src} and \texttt{forth}). The rules and colours are the web site's
  highlighter's, mapped to the terminal's 16 colours; a Forth \texttt{(} comment that runs over lines
  is followed by keeping the state at the start of each row. The keyword tables (mnemonics by
  category, CSR and register names, Forth words) are generated at build time by
  \texttt{tools/vi-syntax.py} from the assembler's own tables, the Forth README and the site's
  highlighter, so the editor, the site and the tools cannot disagree.
\item[The screen's size] vi moves the cursor to row 999, column 999 and asks where it is
  (\texttt{ESC [ 6 n}); the host's terminal and \texttt{term} answer with the bottom right corner.
  Without an answer it is $80\times24$. Ctrl-L asks again.
\item[The text in far memory] The text does not live in the D-LS: it is 256 chunks of 2~KiB in the
  arena, each a run of whole lines, in the order of a chunk table; one chunk at a time is in the D-LS.
  A full chunk is split where a line goes in, an empty one is freed, a small one takes in the next.
  The undo history (a record per change: the old lines and the new) and the register are in the arena
  too. A file of about 500~KiB fits, so the 333~KiB shell source opens; \texttt{fs} writes files of
  128~KiB at most, and \texttt{:w} refuses more before it touches the file.
\item[Drawing] Only the rows that changed, the status line and the cursor are sent, with the cursor
  hidden meanwhile; a scroll draws the text rows again.
\item[Keys and Ctrl-C] The console delivers a key's bytes one at a time, so Esc is told from an
  arrow's \texttt{ESC [ A} by time: vi polls with READs that do not block (which have \texttt{tty}
  read a byte ahead) for 30~ms. vi catches interrupts: Ctrl-C ends insert mode or cancels what is
  being typed, and does not end the editor.
\end{description}

\section{The window server}\label{sec:os-ws}

The graphical interface is a window server, \texttt{ws}, and client programs. \texttt{ws} is the only
process that holds the framebuffer and, while it runs, the keyboard; it draws the desktop and
composites what clients draw. Clients never see the framebuffer. \texttt{ws} is an ordinary program,
started from a shell: it registers as the server named \texttt{ws}, obtains the framebuffer memory and
the controller's registers from \texttt{dev} (GRANT\_FB), makes itself the console's foreground process
in raw mode, and catches interrupts. When it ends, the kernel revokes its framebuffer capability and
the console returns to the shell.

\subsection{Drawing through transfers}
No store can reach the framebuffer, so nothing is drawn in place. The screen is described by a
\term{display list} of items (rectangles, gradients, edges with cut corners, text, blits of window
surfaces, the pointer's cursor). To draw a rectangle of the screen the renderer walks it strip by strip: the items that meet
the strip are collected, an item that covers the strip opaquely drops everything under it, each item
draws its part of each row into a strip buffer in the D-LS, and the strip leaves with one \insn{tput}.
A blit of a window's surface is a \insn{tget} of the source row straight into the strip buffer: no
pixel of a surface passes through a register.

Damage decides which rectangles are drawn. Every frame is: update the model, rebuild the display list,
draw the damaged rectangles, write PRESENT once. While the layout is in motion, frames are paced by
reads of VSYNC; otherwise \texttt{ws} sleeps in RECEIVE with a timeout. Its strip buffers and display
list are a scratch range (Section~\ref{sec:os-switch}).

\begin{center}\small
\begin{tabular}{@{}l r r@{}}
\toprule
\textbf{Update, at 1280 $\times$ 800} & \textbf{Instructions} & \textbf{Pixels}\\\midrule
A full frame: grid, status strip, four tiles and their blits & 3.77 M & 1 024 000\\
A terminal line: render two text rows into the surface and composite them & 161 k & 22 104\\
A pulse of one core's cell in the status strip & 8.7 k & 504\\
\bottomrule
\end{tabular}
\end{center}

\subsection{The desktop}
The screen reads as an instrument panel of the machine and makes its explicitness visible.
\begin{itemize}
\item \textbf{Tiles.} Up to four windows, tiled without overlap: one full; two side by side; three as
  one left and two stacked right; four as a 2 $\times$ 2 grid. Tiles float on a faint grid with 6-pixel
  gaps; each has a luminous one-pixel edge with cut corners, amber when focused. A tag cut into the top
  edge names the window, its process and its core, and a second tag shows the KiB that context switches
  have moved for that process, updating live. When the layout changes, tiles slide to their new places
  over twelve frames.
\item \textbf{The core lattice.} The 28-pixel status strip has one cell per core, lit in the colour of
  the process running on it and dark when idle; the system core and server cores are marked. A context
  switch on a core makes its cell pulse. Beside it are the window list, uptime, the process count,
  context switches per second and the clock.
\item \textbf{Keys} (Alt+key as the window's keyboard sends it, \texttt{ESC [ 27 ; 3 ;} \textit{code}
  \texttt{\textasciitilde}; Alt+arrows \texttt{ESC [ 1 ; 3} \textit{X}). Alt+Enter opens a terminal; Alt+Space opens the
  command palette, a \texttt{launch} line with completion of program names; Alt+1--4 and Alt+arrows move
  the focus; Alt+Q closes the focused window; Alt+Shift+Q quits the window server; Alt+? shows a panel
  of these bindings, which also appears when the splash ends (any key or click dismisses it). Every
  other key is an event for the focused window: Esc and then a key is Esc and that key, whatever the
  timing, and an Alt+key the server does not bind arrives as ESC and the key. Ctrl-C interrupts the
  window's process.
\item \textbf{The pointer} (Section~\ref{sec:pointer}). The device's waiting read belongs to
  \texttt{tty}, alone on its core, which hands the events over with the keys (a READ with mode bit 1).
  \texttt{ws} draws a cursor; a click focuses the tile under it, runs a program from the palette or
  closes it; the wheel moves the palette's selection. The window under the pointer gets presses,
  releases, moves while a button is held, and the wheel.
\item \textbf{Fonts.} An 8 $\times$ 16 body font and a 5 $\times$ 7 small-capitals font for labels,
  both original and in the public domain.
\end{itemize}

\subsection{The client protocol}
A client finds the server with LOOKUP \texttt{"ws"} and sends it the operations of
Table~\ref{tab:os-ws}. A window's \term{surface}, width $\times$ height pixels of 4 bytes in the
framebuffer's format, lives in the server: in the framebuffer window after the screen (2~MiB per
window). Memory capabilities are linear, so one piece of memory cannot be the client's to write and
the server's to read at the same time. The client therefore draws into a \term{band} of its own arena
(its whole arena, as many full-width rows as fit) and lends the band to the server as the attachment
of a DAMAGE request; the server copies the rows into the surface and returns the band with the reply.
While the band is lent the client cannot write it, so the server never reads a half-drawn row: the
ownership rule that prevents data races between cores is also what makes composition tear-free. And
because the surfaces are the server's, it can redraw any window (a slide, the palette's dim) without
asking its client.

\begin{table}[htbp]
\centering\small
\caption{Operations of the window server}\label{tab:os-ws}
\begin{tabularx}{\linewidth}{@{}r l X@{}}
\toprule
\textbf{Op} & \textbf{Name} & \textbf{Request and reply}\\\midrule
40 & CREATE & mode bit 0 = interrupt the caller (SIGNAL) when events wait and it is not in EVENT.
  Data: width, height (ignored: the tile decides), then the title (8 bytes). Reply: the window id, with
  the surface's width and height in the data; $-7$ if there is no room (four windows)\\
41 & DAMAGE & handle = window; data: x, y, width, height, stride; attachment: a band holding those
  width $\times$ height pixels from its offset 0, rows stride bytes apart (0 = width $\times$ 4). The
  server copies them into the window's surface and redraws that rectangle; the band returns with the
  reply. Without an attachment the rectangle is only redrawn\\
42 & TEXT & handle = window; data: x, y in pixels, foreground and background colours ($-1$ = no
  background), then a string (length = 32 + its bytes): drawn by the server with its 8 $\times$ 16 font\\
43 & FILL & handle = window; data: x, y, width, height, colour\\
44 & EVENT & handle = window; mode bit 0 = do not block; length = a timeout in ticks (0 = none). Reply
  data: the kind (1 key, 2 resize, 3 close, 4 focus in, 5 focus out, 6 pointer, 7 wheel), a value
  (the key's bytes, the first in the low byte; for the pointer and the wheel x and y in the surface in
  bits 31:0 and the device's event bytes 4--7 above them), the width (for a key: the number of bytes)
  and the height; $-3$ if none, $-8$ when the timeout passes\\
45 & DESTROY & handle = window\\
46 & TITLE & handle = window; data = the title (8 bytes)\\
\bottomrule
\end{tabularx}
\end{table}

\begin{description}[leftmargin=1em,style=nextline]
\item[\texttt{term}] a terminal emulator. It creates a window, runs \texttt{sh -i} connected by two
  pipes, renders the shell's output into its band, sends key events to the shell's standard input,
  forwards Ctrl-C as SIGNAL, and scrolls back through its history with the wheel.
\item[\texttt{wmandel}] renders the Mandelbrot set into its band, damaging rows as they finish, and
  renders again on a resize.
\item[\texttt{wclock}] a clock that redraws only the band that changes each second.
\end{description}
\texttt{os/lib/gfx.s} is the drawing library both the server and its clients use: surfaces, the strip
renderer, spans with alpha blending, the standard items and text. \texttt{make ws-demo} runs the
server's renderers on the bare machine against a simulated twelve-core system.

\subsection{Terminal output}
Programs write plain bytes to standard output. On the console they go to the host terminal; a
\texttt{term} window interprets this subset and ignores other escape sequences: newline, carriage
return, backspace and tab; \texttt{ESC [ K} (erase to end of line); \texttt{ESC [ n C} and
\texttt{ESC [ n D} (cursor right and left); \texttt{ESC [ 2 J} and \texttt{ESC [ H} (clear, home);
\texttt{ESC [ r ; c H} (position); and \texttt{ESC [ \ldots\ m} with 0, 1, 30--37, 90--97 and 40--47
(colours), \texttt{ESC [ ? 25 l} and \texttt{h} (the cursor off and on while a program draws), and
\texttt{ESC [ 6 n}, which \texttt{term} answers with \texttt{ESC [}\textit{ row}\texttt{ ;}\textit{
column}\texttt{ R} on the program's standard input, as a terminal would type it. The shell's line
editor uses only the first part of this subset; \texttt{vi} uses all of it.

\section{Networking}\label{sec:os-net}

TCP/IP is a server too: \texttt{net} (\texttt{os/servers/net.s}). \texttt{net}~\textit{address}
starts it (\texttt{net 10.0.0.1 \&}); like \texttt{ws}, a program with exactly that name becomes the
server of the lowest empty core, and like \texttt{tty} it is alone there, because it keeps a device read
in flight: the network interface's RX (Section~\ref{sec:netif}). It registers \texttt{"net"}, obtains
the interface's ports from \texttt{dev} (GETNET, op 59, only for the process registered under that
name) and waits in RECEIVE for a request, a packet (the RX read's tag) or its next timer. The first
program that uses the network starts it with address 10.0.0.1 if it is not running, so
\texttt{ping 127.0.0.1} works on a freshly booted machine.

\begin{description}[leftmargin=1em,style=nextline]
\item[IPv4] one interface with a static address (\texttt{ifconfig} changes it), no routing (every
  address is on the link), no fragments. Packets to 127.0.0.0/8 and to the interface's own address
  never reach the device: they wait in a local queue in the server's arena and are handled as if they
  had arrived, so a machine with no link can still talk to itself.
\item[ICMP] echo requests are answered; replies complete PING requests, timed with the timer.
\item[UDP] sockets bound to ports; datagrams of up to 1472 bytes, queued with their sender.
\item[TCP] four sockets in all, each with a 2~KiB receive ring (the advertised window) and a 2~KiB
  send buffer; segments of at most 1024 bytes; in-order receive; the oldest unacknowledged segment is
  sent again after 250~ms, doubling up to 8~s, at most eight times, before the connection is reset;
  listen and accept, connect, close and half-close. Timers are the RECEIVE timeout, computed from the
  clock in \reg{c7}.
\end{description}

\begin{table}[htbp]
\centering\small
\caption{Operations of \texttt{net} (through LOOKUP)}\label{tab:os-net}
\begin{tabularx}{\linewidth}{@{}r l X@{}}
\toprule
\textbf{Op} & \textbf{Name} & \textbf{Request and reply}\\\midrule
1 & READ & TCP: handle = socket; waits for data. Reply: bytes; 0 once the peer has closed; $-1$ after
  a reset\\
2 & WRITE & TCP: the data. Reply: the bytes taken (at least one); waits while the send buffer is full\\
4 & CLOSE & mode bit 0: stop sending only (half-close). The buffered data goes, then FIN\\
50 & STAT & reply: the sockets in use; data: address, the interface's STATUS, packet counters, then
  32 bytes per socket\\
51 & SETIP & handle = the new address\\
52 & PING & handle = address, length = payload bytes, mode = timeout. Reply: the round trip in ticks;
  $-8$\\
53 & CONNECT & handle = address, length = port. Reply: the socket once established; $-1$ refused;
  $-8$ no answer\\
54 & LISTEN & handle = port. Reply: the socket; $-7$ in use\\
55 & ACCEPT & handle = a listening socket. Reply: the next established connection, now the caller's\\
56 & UDP & handle = port (0: any). Reply: the socket\\
57 & SENDTO & data: address, port, payload. Reply: the bytes sent\\
58 & RECVFROM & waits for a datagram. Reply: its bytes; data: sender's address and port, payload\\
60 & RESOLVE & data = a host name. Reply: its first IPv4 address; data: up to four addresses and the
  canonical name; $-4$ no such name, $-8$ no resolver answered\\
61 & SETDNS & handle, length = the two resolvers\\
\bottomrule
\end{tabularx}
\end{table}

A socket belongs to the process that made it. \texttt{proc} sends CLEANUP only to \texttt{fs},
\texttt{tty} and \texttt{pipe}, so \texttt{net} looks at the process list every two seconds while
sockets are open and closes those whose owners have ended. The library's \texttt{net\_init} puts the
endpoint in \reg{c6} (no network call returns a capability, so nothing else lands there), and TCP
sockets are ordinary handles (server register 6): \texttt{sys\_read}, \texttt{sys\_write} and
\texttt{sys\_close} work on them. The programs are \texttt{ping} (until Ctrl-C or \texttt{-c}~$n$),
\texttt{host} (a name's addresses), \texttt{nc} (TCP and UDP, client and listener), \texttt{tcpecho}
(an echo server), \texttt{ifconfig} and \texttt{netstat}; they take host names. Two simulators joined
with \texttt{-N udp:\ldots} are two machines on one link (\texttt{tools/test-net.sh} runs both); with
\texttt{-N nat} the machine is on the real network.

\textbf{DNS.} RESOLVE sends an A query over UDP to port 53 of the first resolver (8.8.8.8 unless the
file \texttt{resolv} names others), with a random ID from a fresh port; the answer must carry the ID
and echo the question; CNAME chains are followed. With no answer in a second the second resolver
(1.1.1.1) is asked, then each once more. Answers are cached for their TTL in the server's arena; dotted
quads and \texttt{localhost} need no query. \texttt{net\_resolve} in the library does it for programs.

\section{Boot}\label{sec:os-boot}

\begin{enumerate}
\item The boot ROM loads the kernel's code and data into every core's Local Stores and the far image
  into the load area, and enters the kernel on every core with \reg{a0} = the core number
  (Section~\ref{sec:boot}).
\item Core 0 becomes the system core. It makes the framebuffer window a revocable region, sets
  \csr{xslot}, and calls the start-up code unit, which is the first unit paged in. Start-up creates the
  kernel shared table as an atomic segment and gives every worker a copy, together with a core-control
  capability for doorbells and the clock capability; makes the arena pool revocable; and reads the directory from the
  disk with the firmware's block routines, the only time the kernel touches the disk.
\item The workers become dispatchers and wait in \insn{tany}.
\item The system core loads \texttt{proc}, \texttt{dev} and \texttt{fs} from the disk itself, since
  \texttt{fs} must exist before anything can be loaded through it, and creates them on cores 1, 2 and 3.
  \texttt{dev} receives the device root; \texttt{fs} asks \texttt{dev} for the disk.
\item \texttt{proc} spawns \texttt{tty} and \texttt{pipe} as servers on the next free cores, then the
  login shell \texttt{sh -l}, and starts a new one whenever it exits.
\end{enumerate}

\texttt{make boot} boots StrataOS on the emulator with 12 cores: the system core, five servers, and six
cores free of servers for threads. The RTL simulator boots it on \texttt{NCORES} cores (8 by default).

\section{Self-hosting}\label{sec:os-as}

A second implementation of the assembler, \texttt{as}, is written in \IS assembly language
(\texttt{os/progs/as.s}). It accepts the \texttt{sasm} language and produces byte-identical STX2 images:
on the bare machine it reads the source from console input and writes the image to console output;
under StrataOS it runs as \texttt{as < prog.s > prog}, a bare-metal program using the redirection
convention of Section~\ref{sec:os-lib}, so its input is limited to 256~KiB and its output to 64~KiB. The
disk carries the source of every example and utility and of \texttt{as} itself (with \texttt{.include}
expanded, since \texttt{as} has no file system), so programs can be rebuilt on the machine, and
\texttt{as} rebuilds itself. \texttt{tools/test-as.sh} checks byte identity for every example, the
kernel, the boot ROM and \texttt{as} itself (at origins 0 and \hex{8000}), a feature test covering every
instruction and directive, a set of inputs both assemblers must reject, and a StrataOS boot that
assembles and runs a program.
